% Lösungen Einheit 5 von 5 – Rechengesetze und Punkt vor Strich (zusammenbau v0.1)
\einheitenkopf{Einheit 5 von 5 · Rechengesetze und Punkt vor Strich}

\erg{26}{a) die Malaufgabe einkreisen \quad b) $5 + 24 = 29$ \quad c) $2{,}5 + 3{,}6 = 6{,}1$ \quad d) $\frac{1}{6} + \frac{4}{12} = \frac{1}{6} + \frac{2}{6} = \frac{3}{6} = \frac{1}{2}$ \quad e) $5 \cdot 4 = 20$ \quad f) $0{,}5 \cdot 2 \cdot 17 = 1 \cdot 17 = 17$ \quad g) $(4{,}5 + 1{,}5) : 4 = 6 : 4 = 1{,}5$ \quad h) erster Term ja, zweiter Term nein (die Oma ist wie ein Kind gerechnet), dritter Term ja}
\erg{27}{a) $(3{,}7 + 6{,}3) \cdot 4 = 10 \cdot 4 = 40$}
\erg{28}{a) Nein, Überschlag $4 \cdot 5 = 20$. Richtig: $3{,}9 \cdot 5{,}1 = 19{,}89$.}
\erg{29}{a) Er hat von links nach rechts gerechnet statt Punkt vor Strich. Richtig: $6 + 4 \cdot 2{,}5 = 6 + 10 = 16$. \quad b) Mal ist die Kurzform für wiederholtes Plus; dieses Päckchen muss erst fertig sein. So ist es vereinbart: Punkt vor Strich. \quad c) $3 \cdot 1{,}20\text{ €} + 2 \cdot 0{,}85\text{ €} = 3{,}60\text{ €} + 1{,}70\text{ €} = 5{,}30\text{ €}$; das Geld reicht nicht.}
